\appendix

%==========================================================================
\section{The three-term identity}\label{app:identity}
%==========================================================================
Let $d_k=(h_k+h_{k+1})/2$ be the distance between neighbouring midpoints. For every quadratic spline on the
grid, with nodal values $U_k=s(m_k)$ and cell values $c_k=s''|_{I_k}$, the interior rows of~\eqref{eq:identity}
read
\begin{equation}\label{eq:interior}
\frac{U_{k+1}-U_k}{d_k}-\frac{U_k-U_{k-1}}{d_{k-1}}
=\frac{h_{k-1}^2}{8d_{k-1}}c_{k-1}+\Big(h_k-\frac{h_k^2}{8d_{k-1}}-\frac{h_k^2}{8d_k}\Big)c_k+\frac{h_{k+1}^2}{8d_k}c_{k+1}.
\end{equation}
The left-hand side is the standard difference of slopes between midpoints, and the right-hand side a
weighted average of $u''$ over three cells. At the left end, with a prescribed value $g_0=s(a)$ or a zero
slope,
\begin{align}
\frac{U_2-U_1}{d_1}-\frac{2(U_1-g_0)}{h_1}&=\Big(\frac{3h_1}{4}-\frac{h_1^2}{8d_1}\Big)c_1+\frac{h_2^2}{8d_1}c_2 &&\text{(value prescribed)},\label{eq:dirrow}\\
\frac{U_2-U_1}{d_1}&=\Big(h_1-\frac{h_1^2}{8d_1}\Big)c_1+\frac{h_2^2}{8d_1}c_2 &&\text{(zero slope)},\label{eq:neurow}
\end{align}
and symmetrically at the right end. The matrix $A$ has positive entries and is strictly diagonally dominant
by columns, so it is always invertible. The collocation problem $s''(m_k)=f(m_k)$ has a unique solution if a
value is prescribed at one or both ends; if the slope vanishes at both ends, it is solvable if and only if
$\sum_kh_kf(m_k)=0$, and unique up to a constant (Result~1 of the Supplementary Material).

Slopes follow by summing the cell values of $u''$. With the left-end slope
$\sigma_0=2(U_1-g_0)/h_1-c_1h_1/4$ (zero for a zero-slope end),
\begin{equation}\label{eq:D1}
s'(m_k)=\sigma_0+\sum_{j<k}h_jc_j+\tfrac12h_kc_k,\qquad s'(b)=\sigma_0+\sum_kh_kc_k,
\end{equation}
which, with $c=D_2U+e$, gives the first-derivative matrix $D_1$ and the end slopes of Algorithm~1.

%==========================================================================
\section{Time stepping for the phase-field equations}\label{app:sav}
%==========================================================================
With $K$ the discrete $-\Delta$, $G=I$ (Allen--Cahn) or $G=K$ (Cahn--Hilliard), $\bar u=2u^n-u^{n-1}$, and a
scalar auxiliary variable $r$~\cite{shen2018scalar}, the stabilized SAV-BDF2 scheme reads
\begin{align}
\frac{3u^{n+1}-4u^n+u^{n-1}}{2\Delta t}&=-G\mu^{n+1},\label{eq:sav1}\\
\mu^{n+1}&=\varepsilon^2Ku^{n+1}+S(u^{n+1}-\bar u)+r^{n+1}\bar b,\label{eq:sav1b}\\
3r^{n+1}-4r^n+r^{n-1}&=\tfrac12\langle\bar b,3u^{n+1}-4u^n+u^{n-1}\rangle_h,\label{eq:sav2}
\end{align}
where $\bar b=f(\bar u)/\sqrt{E_1(\bar u)+C}$ is the scaled nonlinearity, $f(u)=u^3-u$,
$E_1(u)=\langle\frac14(u^2-1)^2,1\rangle_h$, $C>0$, and $S\ge0$ a stabilization constant. On a uniform
grid, for every time step and every $S\ge0$, the scheme has a unique solution, the Cahn--Hilliard scheme
conserves the discrete mass exactly, and a modified energy never increases (Result~5 of the Supplementary
Material). Algorithm~2 lists one step.

\begin{algorithm}[t]\small\hrule\smallskip
\noindent\textbf{Algorithm 2: one SAV-BDF2 time step} (\texttt{ihaar/phasefield.py}).
Precomputed: the transform $\mathcal T$ (DCT-II, or the eigenvectors $V$), the symbols $\hat\lambda$ of
$-K$, $\hat G=1$ (AC) or $\hat\lambda$ (CH), and $\hat S_2=3+2\Delta t\,\hat G(\varepsilon^2\hat\lambda+S)$.
Stored from previous steps: $\hat u^n$, $\hat u^{n-1}$, $u^n$, $u^{n-1}$, $r^n$, $r^{n-1}$.
\step{1} Compute $\bar u=2u^n-u^{n-1}$ and $\bar b=f(\bar u)/\sqrt{E_{1,h}(\bar u)+C}$, and transform, $\hat{\bar b}=\mathcal T\bar b$ (one forward transform).
\step{2} Compute $\hat g=4\hat u^n-\hat u^{n-1}$, $\hat p=(\hat g+2\Delta t\,S\hat G\,\hat{\bar u})/\hat S_2$ and $\hat q=-2\Delta t\,\hat G\hat{\bar b}/\hat S_2$.
\step{3} Update the scalar, $r^{n+1}=\big(4r^n-r^{n-1}+\tfrac12\langle\hat{\bar b},3\hat p-\hat g\rangle\big)/\big(3-\tfrac32\langle\hat{\bar b},\hat q\rangle\big)$, with inner products evaluated in the orthonormal transform domain.
\step{4} Set $\hat u^{n+1}=\hat p+r^{n+1}\hat q$ and $u^{n+1}=\mathcal T^{-1}\hat u^{n+1}$ (one inverse transform).
\smallskip\hrule
\end{algorithm}

